
% !TeX program = xelatex
\documentclass[9pt,a4paper,twocolumn]{extarticle}

% ------------------------------------------------------------
% Page layout and typography
% ------------------------------------------------------------

\usepackage[
  top=18mm,
  bottom=19mm,
  left=17mm,
  right=17mm
]{geometry}

\setlength{\columnsep}{7mm}
\setlength{\parindent}{1.1em}
\setlength{\parskip}{0pt}
\linespread{1.06}

\usepackage{fontspec}
\usepackage{polyglossia}

\setdefaultlanguage{english}
\setotherlanguage{russian}

\setmainfont{Times New Roman}
\setsansfont{Arial}

\setmonofont[Scale=0.86]{Courier New}
\newfontfamily\cyrillicfonttt[
  Script=Default,
  Scale=0.86
]{Courier New}

\usepackage{amsmath,amssymb,mathtools}
\usepackage{booktabs,tabularx,array}
\usepackage{microtype}
\usepackage{csquotes}
\usepackage{titlesec}
\usepackage{caption}
\usepackage{xcolor}
\usepackage{xurl}
\usepackage{hyperref}

\definecolor{linkblue}{HTML}{245A78}

\hypersetup{
  colorlinks=true,
  linkcolor=linkblue,
  citecolor=linkblue,
  urlcolor=linkblue
}

\titleformat
  {\section}
  {\normalfont\bfseries\large}
  {\thesection}
  {0.65em}
  {}

\titleformat
  {\subsection}
  {\normalfont\bfseries\normalsize}
  {\thesubsection}
  {0.6em}
  {}

\titlespacing*{\section}
  {0pt}{1.4ex plus .4ex}{.8ex}

\titlespacing*{\subsection}
  {0pt}{1.1ex plus .3ex}{.45ex}

\captionsetup{
  font=footnotesize,
  labelfont=bf,
  justification=raggedright,
  singlelinecheck=false,
  skip=4pt
}

\setlength{\textfloatsep}{9pt plus 2pt minus 2pt}
\setlength{\intextsep}{8pt plus 2pt minus 2pt}
\setlength{\emergencystretch}{1.5em}

\newcommand{\Exp}{\mathbb{E}}
\newcommand{\DD}{\mathcal{D}}
\newcommand{\Macro}{\textsc{MacroHFT}}
\newcommand{\Earn}{\textsc{EarnHFT}}

\DeclareMathOperator*{\argmax}{arg\,max}

% ------------------------------------------------------------
% Document
% ------------------------------------------------------------

\begin{document}


\renewcommand{\abstractname}{Abstract}

\begin{abstract}
We critically examine \Macro{}, a hierarchical RL system for minute-by-minute cryptocurrency
trading, and partially reproduce its results using the released ETH experts. We study the
assumptions of the trading simulator, assess whether the use of RL is justified, and analyse
trend and volatility specialisation, dynamic mixing of Q-values, and episodic memory.

The simulator executes trades without latency, slippage, or market impact, which limits the
transferability of its results to live trading. Mathematical counterexamples show that a trend
and volatility classification does not determine the optimal trading action. Soft mixing of
Q-values does not, by itself, guarantee faster adaptation. The authors do not separately test
whether volatility specialisation improves on trend-only specialisation, or whether soft mixing
improves on a comparable hard router. They also do not measure response times around stress
events.

On the ETH test set, two trend experts almost always choose the same action: one stays flat and
the other stays long. A fixed expert selected by validation return outperforms every aggregator
we study on this test, although the result depends on the selection criterion. Ridge-Static
($V_1\rightarrow V_2$) achieves $13.22\%$ under the authors' normalisation. This is close to the
reported $12.73\%$ for \Macro{} without memory, but well below the full model's $39.28\%$. We
work with fixed low-level RL experts and do not reproduce the complete system end to end. Our
results do not establish the need for either a regime hierarchy or an additional RL aggregator.
Whether RL adds value over simpler trading models remains an open question.

We regard the memory mechanism as the most promising direction for further research. The
published ablation shows that it raises return and Sharpe ratio and lowers maximum drawdown on
all four assets. It remains to be established whether this improvement comes from making better
use of rare events, smoothing Q-values, or reducing risk. It should also be tested against
existing alternatives in more realistic trading environments.
\end{abstract}


\section{Trading problem}
\label{sec:problem}

\subsection{Trading environment and available actions}

Zong et al.~\cite{zong2024macrohft} consider trading a single cryptocurrency at fixed one-minute
intervals. They formulate the problem as a Markov decision process:
\begin{equation}
  \mathcal{M}
  =(\mathcal{S},\mathcal{A},\mathcal{T},r,\gamma),
  \label{eq:mdp}
\end{equation}
where $\mathcal{S}$ is the state space, $\mathcal{A}$ is the action space, $\mathcal{T}$ is the
transition function, $r$ is the reward, and $\gamma\in[0,1)$ is the discount factor.

The agent has two actions: hold no position (\emph{flat}) or hold a fixed-size long position
(\emph{long}). Let $P_t\in\{0,1\}$ denote the position immediately before the decision, and let
$a_t$ denote the selected target position:
\begin{equation}
  a_t\in\mathcal{A}=\{0,1\},
  \qquad P_{t+1}=a_t.
  \label{eq:action}
\end{equation}
The value $1$ means holding $m$ units of the cryptocurrency, while $0$ means holding no position.
If $a_t\ne P_t$, the simulator executes a fixed-size purchase or sale; otherwise no trade takes
place. In other words, only \emph{taker-only} strategies are considered: the position is changed
by an aggressive order at the price assumed by the simulator. Limit-order-book snapshots are
used as features, but the setting includes no mechanism for executing limit orders.

Let $p_t$ be the closing price of a one-minute interval and let $m$ be the fixed position size.
The authors define the reward as the change in portfolio value after transaction costs:
\begin{equation}
  r_t=m\Bigl[
    a_t(p_{t+1}-p_t)
    -\delta_t\lvert a_t-P_t\rvert
  \Bigr].
  \label{eq:reward}
\end{equation}
Here $\delta_t$ is the cost of switching positions, expressed in price units per unit of the
asset. With a proportional fee $c$, it can be written as $\delta_t=c\,p_t$. The experiments use
$c=0.0002$, or $0.02\%$ of the trade value~\cite{zong2024macrohft}. The simulator charges no fee
when the position remains unchanged.

The space of \emph{trading actions themselves} is much narrower here than in optimal execution
or market-making problems. Every minute, the agent only has to decide whether to hold a fixed
long position or stay in cash. This does not, however, make the optimal policy trivial.

\subsection{State and optimisation objective}

The observations include OHLCV data and five levels of the limit order book. The authors
distinguish features of the current market snapshot from contextual features computed from past
observations. At the high level, they add rolling estimates of trend and volatility to describe
the longer market context (all features are listed in Table~4 of \cite{zong2024macrohft}).

The standard objective of an RL policy is to maximise expected discounted reward:
\begin{equation}
  J(\pi)=\Exp_{\pi}\!\left[
    \sum_{t=0}^{T-1}\gamma^t r_t
  \right],
  \label{eq:objective}
\end{equation}
where $\pi(a\mid s)$ specifies the action distribution in state $s$. In financial terms, the
reward in~\eqref{eq:reward} accumulates trading PnL after fees. The authors also evaluate the
results using return, Sharpe ratio, and maximum drawdown.


\subsection{Limitations of the execution model}
\label{sec:execution_limitations}

Under this setup, the authors optimise a trading policy $\pi$ on a historical price path that is
unaffected by the policy's actions. When the position switches from flat to long or back, the
simulator updates the position, deducts a fee, and continues replaying the same sequence of
market observations. It fills the trade at a specified price regardless of order size, order-book
conditions, or the time between the signal and execution.

These assumptions need not hold in live trading. We consider latency, liquidity, and the effect
of the strategy's own trades on subsequent market behaviour in turn.

\paragraph{Latency and slippage.}

A trading algorithm makes decisions from information received from the exchange. Some time
passes between a market event and the execution of the order sent in response. The data must
reach the strategy server, the algorithm must process it, the order must travel to the exchange,
and the exchange's trading system must accept and execute it. Let the total latency be
\begin{equation}
    \tau =
    \tau_{\mathrm{data}}
    + \tau_{\mathrm{calc}}
    + \tau_{\mathrm{network}}
    + \tau_{\mathrm{exchange}}.
    \label{eq:execution_latency}
\end{equation}

Professional market participants reduce the network component through \emph{colocation}: they
place trading servers in or near the data centre that hosts the exchange infrastructure. This
both shortens and stabilises the transmission time for data and orders. A signal sent from a
remote server crosses additional network hops, making latency longer and less predictable.
Colocation does not, however, remove computation time or the exchange's order-processing time.

The available price may change during $\tau$. For example, an algorithm observes a buy signal at
the quote $p_t^{\mathrm{ask}}$, but its market order is filled at
$p_{t+\tau}^{\mathrm{ask}}$. The difference
\begin{equation}
    \varepsilon_t^{\mathrm{lat}}
    =
    p_{t+\tau}^{\mathrm{ask}}
    -p_t^{\mathrm{ask}}
    \label{eq:latency_slippage}
\end{equation}
measures the buy-side slippage caused by latency. It need not be positive: the price may move in
the trader's favour. Signals, however, often arise from events to which other participants also
react. A late order may then systematically receive a worse price than the one used when the
expected return was assessed.

Latency matters more for strategies that make decisions every second or faster. If a predictable
move lasts only a few hundred milliseconds, even a good forecast may lose its practical value
before the order is filled. \Earn{}~\cite{qin2024earnhft} makes trading decisions at one-second
intervals but does not include such delays in its execution model.

The one-minute interval used by \Macro{}~\cite{zong2024macrohft} makes some strategies less
sensitive to millisecond-scale latency, but does not eliminate the issue. A decision made once a
minute may depend on a fleeting quote change near the end of that minute. Moreover, if an
interval's OHLCV values become available only after it ends, execution at that same interval's
closing price requires separate justification. Assessing the effect of latency on a particular
minute-level policy therefore requires an execution model.

\paragraph{Order-book depth and market impact.}

Even with zero latency, a market order need not be filled entirely at the best available price.

A limit order book lists the prices and quantities at which other participants are willing to
buy or sell the asset. The best ask is the lowest sell price, and the best bid is the highest buy
price. Their difference, the \emph{bid--ask spread}, is one component of the cost of immediate
execution. A market buyer pays the ask, while a market seller receives the bid.

Only a finite quantity is available at the best price level. Suppose a strategy sends a market
order to buy $0.2$ ETH when only $0.05$ ETH is offered at the best price. The exchange fills the
remaining $0.15$ ETH at the next, more expensive levels of the book. The average execution price
is
\begin{equation}
    \bar p_t^{\mathrm{buy}}(q)
    =
    \frac{1}{q}
    \sum_{k=1}^{K}
       p_{t,k}^{\mathrm{ask}}q_{t,k},
    \qquad
    \sum_{k=1}^{K}q_{t,k}=q,
    \label{eq:lob_sweep}
\end{equation}
where $q_{t,k}$ is the quantity actually filled at level $k$. If the best quote alone cannot fill
the order, the average purchase price exceeds the best ask. An analogous effect occurs on the
bid side when selling.

Slippage can therefore arise even without network latency: finite book depth is enough. The
larger a market order is relative to available liquidity, the more price levels it may sweep.

\Macro{} constructs features from five-level order-book snapshots, but its reward function uses
minute closing prices and fixed transaction costs to determine the execution price. In
particular, equation~\eqref{eq:reward} does not account for the cost of sweeping multiple book
levels. The closing price need not equal the best bid or ask when the trade is executed.

Sweeping the book is part of the broader issue of \emph{market impact}: the effect of one's own
trade on the market price. A large market order consumes available liquidity, may shift the best
quotes, and can alter other participants' decisions. Market makers may cancel orders, reduce
displayed size, or revise their prices in anticipation of a further move. As liquidity recovers,
some of the impact may dissipate while some may persist. The distinction between temporary and
permanent impact underlies, among other approaches, the Almgren--Chriss optimal-execution
models~\cite{almgren2001optimal}.

The historical simulator does not model this mechanism. The market component of the next state
is independent of the strategy's action:
\begin{equation}
    \widehat{\mathcal T}
    (s_{t+1}^{\mathrm{mkt}}\mid s_t,a_t)
    =
    \widehat{\mathcal T}
    (s_{t+1}^{\mathrm{mkt}}\mid s_t).
    \label{eq:exogenous_market}
\end{equation}
In a real market, the distribution of future quotes may depend on the direction and size of the
strategy's own order, as well as on other participants' response to it. Historical data reveal
market dynamics without our intervention, but do not necessarily describe those dynamics after
we intervene.

This distinction becomes important as capital grows. In~\eqref{eq:reward}, monetary profit scales
linearly with the position size $m$. In reality, increasing $m$ may worsen the execution price
more sharply and alter subsequent liquidity, breaking this linear scaling. A strategy that is
profitable at small size may lose money when substantially more capital is allocated to it.

For a fixed order size of $0.2$ ETH in a liquid market, the impact of a single trade may be
small. The simulator's results nevertheless reveal neither the size of that effect nor the
capital scale at which the policy retains its edge.

\paragraph{Restriction of the strategy space.}

Under~\eqref{eq:action}, the agent can change its position only through an aggressive trade. This
form of execution reaches the desired position relatively quickly, but requires accepting the
market price and paying a taker fee.

A market participant can instead place a limit order to buy at the bid or sell at the ask and
wait for an incoming order. It then \emph{provides liquidity} as a maker rather than consuming it
as a taker. If a purchase and a subsequent sale are both filled favourably, the participant can
earn the bid--ask spread instead of paying it twice. Some exchanges and fee tiers also offer
\emph{maker rebates}, payments or discounts for supplying liquidity.

Passive execution is not a cost-free substitute for a market order. A limit order may never be
filled: the price can move in the anticipated direction without reaching it. Earlier orders may
be ahead in the queue at the same price, so merely touching that price does not guarantee a
fill. Finally, a passive order is particularly likely to execute when informed traders are
willing to trade against it. This creates \emph{adverse selection}; for example, a bid-side buy
may fill just before the price falls further.

High-frequency liquidity-provision strategies must therefore manage queue position, order
lifetime and cancellation, non-execution risk, and accumulated inventory. At the same time, for
strategies with small expected price moves, savings on the spread and fees can determine whether
profit is possible at all. The taker-only restriction excludes such choices, including placing
several limit orders at different price levels and working a large order in smaller pieces.

At a fee rate of $c=0.0002$, a purchase followed by a sale already costs about four basis points
of traded value, even before the spread and slippage. If the available predictive signal cannot
cover these costs, no profitable active taker policy may exist. Passive execution strategies can
remain potentially profitable under the same market conditions. Thus, solving
\[
    \pi^*_{\mathrm{taker}}
    =
    \argmax_{\pi\in\Pi_{\mathrm{taker}}}J(\pi)
\]
does not guarantee finding the best strategy in the broader space of available trading actions
$\Pi_{\mathrm{all}}$.

\paragraph{Transferring results beyond the simulator.}

These simplifications can change not only the absolute returns of the strategies but also their
relative ranking. Let $\widehat J(\pi)$ denote the result of policy $\pi$ in the historical
simulator, and let $J_{\mathrm{real}}(\pi)$ denote its expected result under live execution with
the same capital.

The inequality
\[
    \widehat J(\pi_A)>
    \widehat J(\pi_B)
\]
does not, in general, imply
\[
    J_{\mathrm{real}}(\pi_A)>
    J_{\mathrm{real}}(\pi_B).
\]

For example, suppose two policies earn $10$ and $8$ basis points of capital over the same
historical period after the fees included in the simulator. The first makes $20$ trades and the
second makes two. Suppose live execution imposes an additional loss of one basis point of
capital per trade. Then
\begin{equation}
    \begin{aligned}
      \widehat J(\pi_A)&=10,
      &J_{\mathrm{real}}(\pi_A)&=-10,\\
      \widehat J(\pi_B)&=8,
      &J_{\mathrm{real}}(\pi_B)&=6.
    \end{aligned}
    \label{eq:policy_ranking_reversal}
\end{equation}
The numbers are illustrative. The example shows that a policy with an advantage in the simulator
may trail a more cautious policy once execution is taken into account. At larger order sizes,
the difference may be greater still because the additional cost need not remain constant per
trade.

A return improvement from adding an RL component establishes its usefulness primarily within
the trading environment being modelled. Whether the effect survives latency, limited liquidity,
and changes in capital remains a separate empirical question.

\paragraph{Directions for further research.}

One approach is to enrich the mathematical model of execution. Instead of a fixed fee, one can
use a cost function that depends on order size, spread, available book depth, and latency:
\begin{equation}
    C_t
    =
    C(s_t,q_t,\tau_t),
    \qquad
    q_t=m|a_t-P_t|.
    \label{eq:execution_cost_model}
\end{equation}
A richer model can include temporary and permanent price impact, the probability of a partial
fill, and liquidity recovery. Its parameters can be estimated from order-book and trade data.
This approach makes it possible to study sensitivity to specific execution characteristics at a
relatively modest computational cost. Its impact function must, however, be chosen in advance,
and other participants' response may be poorly captured by a simple parametric relationship.

A second approach is \emph{agent-based modelling} (ABM). In such an environment, a virtual
exchange processes individual orders, cancellations, and trades, while a population of trading
agents generates order flow. The policy under study interacts with them through the same
execution mechanism, so its actions can change both the book and future observations.

ABIDES~\cite{byrd2020abides} is one example of this infrastructure: a discrete-event market
simulator that models the order book, participant interactions, and network latency. Events can
carry nanosecond-resolution timestamps, allowing the arrival order of messages to be reproduced
at scales far below one minute. Nanosecond precision in the clock does not by itself make the
participants' behaviour realistic.

The main challenge for ABM is calibration. A useful environment must reproduce not only the
return distribution but also order flow, book depth, order lifetimes, the rate of liquidity
recovery, and participants' reactions to large trades. Such an environment is much more
computationally demanding than replaying historical minute observations. Even a well-calibrated
agent-based market remains a model rather than an exact copy of a live exchange.

Both directions permit a more concrete test of the practical value of trading RL architectures:
does policy $\pi$ retain its advantage as latency $\tau$, position size $m$, and the execution
mechanism change? When \Macro{} is compared with simpler policies, such a test may alter the
conclusions drawn from a single historical backtest even if the trained models themselves remain
unchanged.

\section{The MacroHFT architecture: motivation and evaluation}
\label{sec:method}

\Macro{}~\cite{zong2024macrohft} continues a line of research on RL for trading that includes,
among others, \cite{deng2016deep,theate2021trading,zhu2022perturbation}. Its immediate predecessor
is \Earn{}~\cite{qin2024earnhft}: the system trains a pool of low-level policies, while a separate
RL router selects the agent that controls trading. Low-level decisions are made every second and
the agent is selected once a minute.

In \Macro{}, both levels operate at one-minute intervals. The authors construct a pool of six
agents trained on trend and volatility subsets, and the high-level network mixes their Q-values
rather than selecting a single expert. We next examine which problems these changes are meant to
solve and which experiments can test the proposed explanations.

\subsection{Adding volatility to the market-regime classification}
\label{sec:regime_decomposition}

\paragraph{Motivation and architectural choice.}

Regime specialisation rests on the premise that different trading policies are useful in
different parts of the market. \Macro{} refines this idea: classifying regimes by trend alone is
insufficient because markets moving in a similar direction may have different volatility. The
authors argue that ignoring this distinction weakens both specialisation and risk
management~\cite{zong2024macrohft}.

They divide the historical data into blocks of length $L$. For each block, they compute the slope
of a smoothed price series and the volatility of returns. Thresholds determined on the train set
define three trend categories and three volatility categories, producing the subsets
\begin{equation}
\begin{aligned}
  \mathcal D^{\mathrm{tr}}_1,\;
  \mathcal D^{\mathrm{tr}}_2,\;
  \mathcal D^{\mathrm{tr}}_3,\\
  \mathcal D^{\mathrm{vol}}_1,\;
  \mathcal D^{\mathrm{vol}}_2,\;
  \mathcal D^{\mathrm{vol}}_3.
\end{aligned}
\label{eq:regime_subsets}
\end{equation}
The authors train a separate dueling DDQN network on each subset. These are two three-category
classifications, not the nine intersections of trend and volatility. They set $L=4320$ minutes
for ETH, DOT, and LTC, and $L=360$ minutes for BTC~\cite{zong2024macrohft}.

\paragraph{What needs to be established.}

A difference in market statistics does not by itself explain why separate policies should be
trained on the corresponding subsets. What matters to a trading decision is the conditional
consequence of each action:
\begin{equation}
  \Delta^*(s)
  =
  Q^*(s,1)-Q^*(s,0).
  \label{eq:optimal_action_difference}
\end{equation}
With the binary long/flat choice, the sign of $\Delta^*(s)$ determines the action and its
magnitude measures the advantage of one action over the other. A trend--volatility partition is
useful for specialisation if it makes such differences easier to estimate. The authors do not
separately test whether their chosen classification improves policy learning.

A simple counterexample illustrates the problem. Consider two histories of five consecutive
prices:
\begin{equation}
\begin{aligned}
  H_A &= p(1,k,k^2,k,1),\\
  H_B &= p(1,k^{-1},k^{-2},k^{-1},1),
  \qquad k>1.
\end{aligned}
\label{eq:matched_histories}
\end{equation}
Both histories end at their initial price $p$. By symmetry, the slope from regressing price on
time is zero in both cases. The simple returns are
\[
  (k-1,\;k-1,\;k^{-1}-1,\;k^{-1}-1)
\]
for $H_A$, and the same sequence in reverse order for $H_B$. The histories therefore have the
same mean return and standard deviation.

Their most recent moves point in opposite directions: the price is falling in $H_A$ and rising
in $H_B$. Suppose the conditional dynamics of the next return $u_{t+1}$ exhibit positive
autocorrelation:
\begin{equation}
  \Exp[u_{t+1}\mid H_t]=\rho u_t,
  \qquad \rho>0.
  \label{eq:momentum_counterexample}
\end{equation}
Starting from flat, the agent then remains in cash after $H_A$ but buys after $H_B$, provided the
expected move covers the entry cost. Here, identical linear trend and volatility statistics
correspond to different preferred actions.

This example uses a linear trend without preliminary filtering. It demonstrates the information
lost by summary statistics: the order and conditional predictability of moves are not determined
by their mean and dispersion. Smoothing or changing the window length is no substitute for
testing whether the chosen classification preserves distinctions relevant to trading decisions.

The converse is also possible. Two regimes may differ in return variance but have the same
conditional expectation. A one-period policy that maximises expected PnL under identical costs
may then choose the same action in both regimes. A difference in volatility alone does not
require different neural networks or different policies.

\paragraph{The missing control.}

Comparing \Macro{} with \Earn{} does not isolate the effect of adding volatility: the expert
pool, the hierarchy's timing, the aggregation rule, and the training procedure all change at
once. Table~3 of~\cite{zong2024macrohft} ablates other components, but provides no direct
trend-only versus trend-plus-volatility comparison.

Moreover, moving from three experts to six increases the number of trained models. Even a gain
from this system would therefore not establish that volatility specialisation itself is useful.
A control should use the same number of experts---for example, several independently trained
agents for each trend subset, or an alternative partition of comparable size.

Testing the idea requires separating three effects: adding volatility information, changing the
training sample, and increasing the number of models. Without such a comparison, the extra axis
of specialisation remains an engineering hypothesis whose contribution is unknown.

\subsection{The policies produced by separate training}
\label{sec:expert_specialization}

\paragraph{From training regime to expert behaviour.}

After segmentation, the authors train six policies $\pi_1,\ldots,\pi_6$. Each receives its own
subset but uses the same action space $\mathcal A=\{0,1\}$. A regime determines which examples
the network sees in training; it does not restrict where the trained network can be applied.

This raises two distinct questions. Does an agent gain a comparative advantage in its assigned
regime? And can it recognise when conditions have changed and its previous action is no longer
profitable? The first property does not guarantee the second.

The released ETH trend agents make this distinction particularly clear in our reproduction.
Table~\ref{tab:conditional_exposure} reports the fraction of time spent long within retrospectively
labelled Bear, Neutral, and Bull segments of the test period.

\begin{table}[t]
  \centering
  \caption{Exposure of trend experts on the ETH test set: percentage of minutes spent long.
  Each expert trades its own sequence of positions.}
  \label{tab:conditional_exposure}
  \footnotesize
  \renewcommand{\arraystretch}{1.15}
  \begin{tabular}{@{}lrrr@{}}
    \toprule
    Expert & Bear & Neutral & Bull\\
    \midrule
    \texttt{slope/1} & 9.66 & 5.14 & 11.04\\
    \texttt{slope/2} & 26.44 & 26.71 & 49.96\\
    \texttt{slope/3} & 99.88 & 99.73 & 99.64\\
    \bottomrule
  \end{tabular}
\end{table}

The \texttt{slope/3} agent holds ETH almost continuously, including during bear segments. By
contrast, \texttt{slope/1} stays mostly in cash regardless of the trend. The same broad pattern
appears on validation: the agents' exposures are 99.56\% and 5.84\%, respectively. It is not
confined to one unusual part of the test period.

Such policies can indeed perform well in their own regimes. An almost constant long position
earns money in a rising market, while an almost constant flat position avoids losses in a falling
one. This outcome does not require an expert to recognise a regime change; a stable preference
for one position is enough.

Our conditional-PnL decomposition is consistent with this interpretation: about 85--87\% of the
measured contrast in relative specialisation between the two extreme trend agents is associated
with the market-exposure component. This is a share of the specialisation measure, not of the
agent's total profit.

On the trajectories we study, part of the pool operates mainly as a set of directional long/flat
preferences. The high level must recognise when each preference should be trusted. Merely
training a bull expert on bull-market data does not create a mechanism for exiting when the
market turns bearish.

\paragraph{Does this require separate RL networks?}

A complex neural network is unnecessary for an almost constant position. Yet a constant action
does not prove that the network's numerical Q-values are useless: even if the sign of the long
advantage never changes, its magnitude may carry information useful to the hyper-agent.

A targeted ablation is therefore needed: does system performance survive if the extreme
experts' recommendations are replaced by constant flat/long policies, or if the magnitude of
their Q-advantages is made independent of state? This would separate the value of continuous
estimates from the value of position diversity itself.

The other experts' results likewise do not show that partitioning the data requires a hierarchy.
\texttt{slope/2} achieves $23.87\%$ over the full ETH test set under the authors' return
normalisation and outperforms every aggregator we consider. A useful policy trained on one
subset can therefore be applied far beyond its original category.

The reason for this result requires a separate test. The subset may have made training easier,
the resulting policy may have transferred better to the new period, or the outcome may depend on
the particular training run. Testing whether separate training is necessary requires an agent of
the same architecture trained on the entire market, together with multiple random initialisations.
A single policy with regime features or balanced training examples is a direct alternative to six
independent networks.

\subsection{Replacing expert selection with weighted Q-values}
\label{sec:q_aggregation}

\paragraph{The proposed solution.}

In the abstract of \Macro{}, the authors describe a single agent's decisions as
\foreignlanguage{english}{\enquote{one-sided and highly biased}}. They associate this property
with possible losses in extreme market conditions~\cite{zong2024macrohft} and propose considering
several experts' estimates at once.

Let $s_t$ denote the available state, including the current position. Each network receives the
part of this state intended for it. Hard routing takes the form
\begin{equation}
\begin{aligned}
  i_t &= g(s_t),\\
  \pi_{\mathrm{hard}}(s_t)
      &=\argmax_{a\in\{0,1\}}
        Q_{i_t}(s_t,a).
\end{aligned}
\label{eq:hard_routing}
\end{equation}
The router $g$ selects one expert, after which the system executes that expert's preferred action.

In \Macro{}, the high-level network produces softmax weights rather than an index:
\begin{equation}
  w_i(s_t)=
  \frac{\exp f_i(s_t)}
       {\sum_{j=1}^{6}\exp f_j(s_t)}.
  \label{eq:mixing_weights}
\end{equation}
The aggregate Q-value is
\begin{equation}
  Q_{\mathrm{mix}}(s_t,a)
  =
  \sum_{i=1}^{6}
  w_i(s_t)Q_i(s_t,a).
  \label{eq:mixed_q}
\end{equation}
With two actions, only the differences
\begin{equation}
  d_i(s_t)
  =
  Q_i(s_t,1)-Q_i(s_t,0),
  \label{eq:action_advantages}
\end{equation}
are needed, so the policy becomes
\begin{equation}
  \pi_{\mathrm{mix}}(s_t)
  =
  \mathbf 1\left\{
    \sum_{i=1}^{6}
    w_i(s_t)d_i(s_t)>0
  \right\}.
  \label{eq:mixed_policy}
\end{equation}
The policy chooses flat when the sum is exactly zero.

\paragraph{What actually changes.}

Equation~\eqref{eq:mixed_policy} does not diversify capital across six strategies. Its output is
still a single position: either $0$ or $m$ units of the asset. Soft mixing changes the rule for
selecting that position, but does not create an intermediate exposure.

Moreover, in this binary problem, every action taken by the mixture agrees with the action of at
least one expert:
\begin{equation}
  \pi_{\mathrm{mix}}(s)
  \in
  \{\pi_1(s),\ldots,\pi_6(s)\}.
  \label{eq:binary_action_support}
\end{equation}
If the weighted sum of $d_i(s)$ is positive, at least one $d_i(s)>0$. If the sum is non-positive,
at least one expert allows the flat action. Given the same information, a sufficiently expressive
hard router can select an expert that takes the same action.

This observation does not imply that the two mechanisms are equally easy to learn. It does rule
out the explanation that hard routing is inherently unable to produce the required defensive
action whereas soft mixing expands the set of available choices. Any advantage must instead come
from how the mechanisms learn and use the experts' estimates.

Numerical mixing has a cost of its own. An expert's influence is determined by the product $w_i
d_i$, not by its weight alone. For example,
\[
  w_1=0.9,\quad d_1=0.1,
  \qquad
  w_2=0.1,\quad d_2=-2
\]
give a total advantage of $-0.11$ and hence the flat action. The expert with a 10\% weight
determines the action in opposition to the expert with a 90\% weight.

A Q-advantage estimates a difference in future rewards, not the probability that a recommendation
is correct. A large absolute value may reflect either an economically meaningful difference
between actions or an estimation error. A plot of weights alone therefore does not show which
expert determined the action; at minimum, the products $w_i d_i$ and their sum are needed.

Positive rescaling of a Q-function preserves an individual expert's policy but can change the
ensemble's action at fixed weights. The quality and calibration of numerical estimates become a
separate requirement, one that a simple vote over the signs of $d_i$ does not have. The authors
do not investigate this issue by comparing raw Q-values, their normalised versions, and the
policies' own recommendations.

\subsection{The claimed faster adaptation is not measured}
\label{sec:adaptation_claim}

\paragraph{The original problem.}

In the introduction to \Macro{}, the authors cite individual policies' delayed response to sudden
fluctuations as one reason for changing the architecture. In the results section, they attribute
to the new architecture an \foreignlanguage{english}{\enquote{ability to adjust quickly to sudden
market changes}}~\cite{zong2024macrohft}. Does the proposed system reduce the delay before taking
an appropriate trading action after a sharp market change?

Aggregate return over several months cannot answer this question. A policy may identify profitable
periods in quiet markets more successfully while responding just as slowly to crashes. The
opposite is also possible: earlier exits can reduce crash losses but lower returns through false
alarms.

\paragraph{Mixing does not eliminate recognition delay.}

Both high-level mechanisms can update their decision at the next one-minute step. If the market
features do not reveal a change in time, learned weights retain the same problem. Nor is it enough
for the pool to contain a suitable expert: that expert's recommendation must exert sufficient
influence on the final action.

This is clear even with two experts. Suppose one prefers long and the other flat:
\[
  d_{\mathrm{long}}=u>0,
  \qquad
  d_{\mathrm{flat}}=-v<0.
\]
Let $b_t$ be the weight of the second expert. Selecting the expert with the largest weight leads
to a switch to flat when $b_t>1/2$. Under Q-value mixing, the exit condition is
\begin{equation}
\begin{aligned}
  (1-b_t)u-b_t v &\leq 0,\\
  b_t &\geq \frac{u}{u+v}.
\end{aligned}
\label{eq:adaptation_threshold}
\end{equation}

If $v=u/4$, exiting requires an 80\% weight on the flat expert. If $v=4u$, 20\% is enough. For
the same path of changing weights, the mixture can therefore switch either later or earlier than
hard selection. Faster response does not follow from the mixing operation itself; it depends on
the experts' estimates and on the learned rule for using them.

This counterexample does not describe the measured behaviour of \Macro{}. It refutes the general
explanation that ``several weighted experts respond faster than a single selected expert.''

\paragraph{What the paper presents.}

Figure~4 contains individual examples of trading decisions, including an exit during a fall in
LTC. For the same event, however, the authors do not show \Earn{}'s response, analyse a
systematically selected sample of stress episodes, or report a distribution of exit delays.

Figure~5 shows expert weights averaged over 60 minutes. Such a plot cannot reveal whether the
weights changed before a sharp move, one minute afterwards, or much later within the same hour.
Moreover, weights without Q-advantages do not determine when the trading position changed.

Even the aggregate risk metrics show no unconditional advantage for the new system. In Table~1,
maximum drawdown on DOT is 15.89\% for \Macro{} versus 5.80\% for \Earn{}; on LTC it is 14.24\%
versus 9.63\%~\cite{zong2024macrohft}. These figures do not measure adaptation delay, and they
cannot support a claim that the new architecture systematically provides better protection from
large losses.

The authors do not test the claimed link between the architecture and rapid adaptation. They
report final trading results and behavioural illustrations, but do not measure the advantage
used to motivate the change in policy-selection mechanism.

\paragraph{An experiment that would test the claim.}

The benchmark set already contains a hard-routing baseline: \Earn{} is one of the algorithms
being compared. The missing control is a learned hard router and a learned soft mixer operating
on the same pool of fixed experts, with identical inputs and trading frequency and comparable
training.

Such an experiment should identify stress events using a rule fixed in advance. It should then
measure the delay before a position change, losses in the event window, and the rate of false
defensive exits outside those windows. The initial position and the opportunity to reduce risk
must be handled explicitly: an agent that was already in cash before an event did not make a fast
exit.

An additional control would freeze the mixture weights before an event while allowing the
experts' Q-values to change. Comparing this version with dynamic weights would show whether
reweighting itself helps, or whether the low-level networks provide the defensive action. The
system's total PnL cannot separate these mechanisms.

\subsection{What simple aggregation alternatives show}
\label{sec:aggregation_experiments}

\paragraph{The substitution being tested.}

In our experiments, we hold fixed the six released ETH experts, the authors' features, the
position size, and the fee. We change only the rule that turns the experts' estimates into a
trading action. Each strategy passes its own previous position to the experts.

In addition to a uniform mixture, majority vote, and normalised Q-advantages, we test a single
fixed expert, a tabular hard router, and a supervised Ridge controller. The tabular router is
trained on the nine combinations of trend and volatility levels and selects one of the six
experts. It is a simple learned alternative, not a reproduction of \Earn{}'s RL router.

Nor can Ridge be described merely as replacing softmax weights with fixed coefficients. In the
notebook, Ridge predicts the next one-minute return:
\begin{equation}
  \widehat{\mu}_t
  =
  b+\sum_{i=1}^{6}
  \beta_i\frac{d_i(s_t)}{\sigma_i},
  \label{eq:ridge_forecast}
\end{equation}
where the scales $\sigma_i$ are estimated on the meta-model's training segment. We then apply a
fee-aware decision rule:
\begin{equation}
  \pi_{\mathrm{Ridge}}(s_t)
  =
  \argmax_{a\in\{0,1\}}
  \left\{
    a\widehat{\mu}_t
    -c|a-P_t|
  \right\}.
  \label{eq:ridge_policy}
\end{equation}
From flat, the rule enters when $\widehat{\mu}_t>c$; from long, it exits when
$\widehat{\mu}_t\leq-c$. This is a distinct and simple way to manage a position using the outputs
of RL experts.

We train the main Ridge variant on the first half of validation ($V_1$) and select its
regularisation on the second half ($V_2$). The contextual version adds interactions between the
expert signals and rolling trend and volatility. The low-level policies remain unchanged in both
cases.

\begin{table*}[t]
  \centering
  \caption{Selected results for alternative ways of using the six ETH experts on the test set
  from June to October 2023. MTM Return, Sharpe, and MDD use a common initial capital;
  Author-style Return uses \texttt{required\_money} as its denominator. Exposure is the
  percentage of minutes spent long.}
  \label{tab:aggregation_results}
  \footnotesize
  \renewcommand{\arraystretch}{1.15}
  \begin{tabularx}{\textwidth}{@{}Xrrrrr@{}}
    \toprule
    Method
    & MTM, \%
    & Author-style, \%
    & Sharpe
    & MDD, \%
    & Exposure, \%\\
    \midrule
    Uniform Q
      & 5.72 & 5.14 & 0.773 & 12.92 & 29.64\\
    Normalized Q
      & -4.04 & -3.53 & -0.333 & 18.07 & 34.31\\
    Majority
      & 9.08 & 9.31 & 1.128 & 9.65 & 35.38\\
    Tabular hard router
      & 10.04 & 9.94 & 0.915 & 18.17 & 82.93\\
    Ridge-Static, $V_1\rightarrow V_2$
      & 13.36 & 13.22 & 1.414 & 12.96 & 46.80\\
    Ridge-Contextual, $V_1\rightarrow V_2$
      & 6.88 & 7.22 & 0.992 & 9.37 & 24.67\\
    Fixed expert:
    selected by validation return
      & 23.28 & 23.87 & 2.424 & 9.54 & 32.09\\
    Fixed expert:
    selected by validation Sharpe
      & -22.47 & -18.74 & -3.116 & 25.28 & 22.80\\
    \bottomrule
  \end{tabularx}
\end{table*}

\paragraph{A result that needs explaining.}

On this test, the best variant we consider uses the \texttt{slope/2} expert selected by
validation return and never switches experts at all. Its MTM Return is 23.28\%, compared with
9.08\% for Majority. Exposure is slightly lower (32.09\% versus 35.38\%), while maximum drawdown
is almost identical (9.54\% versus 9.65\%). The gap cannot be explained simply by spending more
time long.

At the same time, the result depends on the expert-selection procedure. The original selection
by validation Sharpe chose \texttt{vol/2} and produced a negative test return. The observation
therefore establishes a narrower fact, not the general reliability of fixed selection: among the
released policies, one expert performs better in this experiment than all of our combinations.
Such a baseline is necessary when evaluating the incremental value of aggregation.

Static Ridge achieves 13.22\% under the authors' normalisation. In Table~3, the authors report
12.73\% for \Macro{} without memory and 39.28\% for the full model~\cite{zong2024macrohft}. A
simple supervised alternative thus reaches a result numerically comparable to the published
no-memory configuration. This is a reason to include it in a direct comparison, rather than
attribute the full system's entire result to a learned voting mechanism.

The remaining tests do not support the idea that any simplification will do. Normalising the
Q-advantages worsens the result. The tabular hard router barely improves on Majority in return
but substantially increases drawdown and exposure. Contextual Ridge trails Static Ridge in
return. Our evidence supports neither a universal advantage for hard routing nor the claim that
mis-scaled Q-values explain every problem.

\subsection{Episodic memory in hyper-agent training}
\label{sec:memory}

\subsubsection{Motivation}

After training specialised experts and introducing a hyper-agent, the stability of the latter's
training remains a concern. In Section~4.3 of \Macro{}~\cite{zong2024macrohft}, the authors
identify two causes. The first is high reward variance among similar market states. The second is
rare, sharp fluctuations that standard experience replay is presumed not to use effectively
enough.

Experience replay already stores past transitions and reuses them during training. Uniform
sampling, however, reproduces the frequencies in the observations: ordinary states appear in
training batches far more often than rare ones. If a few episodes of large losses determine the
quality of the entire strategy, this allocation of updates may be undesirable.

The authors propose augmenting training with episodic memory. For the current state, the method
retrieves similar transitions and uses their stored value estimates as an additional training
signal. The expectation is that this will align the Q-values of nearby states and help produce a
policy that responds more quickly to sharp market changes.

This approach extends an existing line of research. Episodic Memory Deep Q-Networks already used
memory to provide additional training for a Q-network~\cite{lin2018episodic}, while Neural
Episodic Control combined a learned state representation with a table of value
estimates~\cite{pritzel2017neural}. \Macro{} cites both studies. The research question is which
problem of the trading hyper-agent is solved by this particular memory design, and what explains
its advantage.

\subsubsection{Memory design and operation}

Let $Q_\theta(s,a)$ denote the hyper-agent's Q-function, where $\theta$ are the parameters of the
trainable high-level network. As before, this network constructs its Q-function by weighting the
low-level experts' estimates. The memory is stored separately from the network parameters:
\begin{equation}
  \mathcal M_{\mathrm{mem}}=(K,E,V).
  \label{eq:memory_structure}
\end{equation}
$K$ contains retrieval keys, $E$ contains states and actions, and $V$ contains stored action-value
estimates.

\paragraph{Adding a transition.}

For an observed transition $(s_j,a_j,r_j,s'_j)$, the encoder computes the key
\begin{equation}
  h_j=\psi_{\theta_j}(s_j),
  \label{eq:memory_key}
\end{equation}
where $\psi_{\theta_j}$ is the hyper-agent's state encoder at the time the entry is created. For a
non-terminal transition, the stored value is
\begin{equation}
  v_j
  =
  r_j+\gamma\max_{a'}Q_{\theta_j}(s'_j,a').
  \label{eq:memory_stored_value}
\end{equation}
In a terminal state, the code sets the future reward to zero.

$v_j$ is not the observed profit from the rest of the trade or the completed episode. Only $r_j$
is observed directly; the network computes the remainder from its own Q-value. This is a one-step
bootstrapped target stored for later training.

The table has finite capacity. When it is full, the code removes the oldest entry under a FIFO
rule. The memory thus retains recent transitions rather than a permanent archive of the most
dangerous market events.

\paragraph{Retrieving similar states.}

For state $s$, the code computes the current key $h=\psi_\theta(s)$ and then selects the $k$
stored keys closest in Euclidean distance. Let $\mathcal N_k(h)$ be the set of their indices.

The published implementation uses the kernel
\begin{equation}
  \kappa(h,h_j)
  =
  \frac{1}{\|h-h_j\|_2^2+\varepsilon},
  \qquad \varepsilon>0.
  \label{eq:memory_kernel}
\end{equation}
Closer entries receive more weight. Among the retrieved neighbours, the implementation then uses
only entries with the same action:
\begin{equation}
  \omega_j(s,a)
  =
  \frac{
    \kappa(h,h_j)\mathbf 1\{a_j=a\}
  }{
    \displaystyle
    \sum_{\ell\in\mathcal N_k(h)}
    \kappa(h,h_\ell)\mathbf 1\{a_\ell=a\}
  }.
  \label{eq:memory_weights}
\end{equation}
The auxiliary estimate is
\begin{equation}
  Q_M(s,a)
  =
  \sum_{j\in\mathcal N_k(h)}
  \omega_j(s,a)v_j.
  \label{eq:memory_estimate}
\end{equation}

Equations~\eqref{eq:memory_kernel}--\eqref{eq:memory_weights} match the published
code~\cite{macrohftcode}. Equation~(7) in the paper uses the distance itself rather than its
inverse. This difference changes the neighbours' relative contributions; we use the implemented
version here. Until the table is full, or if none of the retrieved neighbours has the relevant
action, the code substitutes the ordinary one-step estimate for $Q_M$.

\paragraph{Policy training and execution.}

The loss function includes the retrieved target:
\begin{equation}
  \begin{aligned}
    \mathcal L(\theta)
    &=
    \mathcal L_0(\theta)
    +\lambda_M\mathcal L_M(\theta),\\
    \mathcal L_M(\theta)
    &=
    \mathbb E\left[
      \bigl(Q_\theta(s,a)-Q_M(s,a)\bigr)^2
    \right].
  \end{aligned}
  \label{eq:memory_loss}
\end{equation}
Here $\lambda_M\geq0$ controls the strength of the auxiliary training, while $\mathcal L_0$
collects the original method's other terms, including the TD error and an auxiliary DP target.
The implementation computes $Q_M$ when collecting a transition and stores it with that transition
in the ordinary replay buffer.

During test-time trading, the model performs no memory lookup: the trained expert and hyper-agent
networks determine the action~\cite{macrohftcode}. Memory affects the outcome through the
parameters learned during training. The explanation that ``the agent saw a crash, recalled a
similar case, and exited the position'' therefore does not describe this model's execution.

The extra training may still improve responses by producing a more suitable policy. But the
speed at which transitions are learned during training and the speed at which the policy exits a
position in a new market are different properties. The former does not establish the latter.

\subsubsection{Published results}

The authors test the value of memory with the \emph{w/o-MEM} ablation. For memory, they directly
compare the full system with a version lacking that component; no equivalent comparison is
provided for volatility specialisation or soft aggregation.

Table~\ref{tab:memory_ablation} reproduces the results from Table~3 of
\cite{zong2024macrohft}. We did not retrain the hyper-agent as part of this review; the table
contains the published results, not our own memory experiment.

\begin{table*}[t]
  \centering
  \caption{Memory ablation in \Macro{}, from Table~3 of~\cite{zong2024macrohft}.
  TR is return, ASR is the annualised Sharpe ratio, and MDD is maximum drawdown.
  All values are taken from the paper.}
  \label{tab:memory_ablation}
  \footnotesize
  \renewcommand{\arraystretch}{1.15}
  \begin{tabular}{@{}lrrrrrr@{}}
    \toprule
    & \multicolumn{3}{c}{Without memory}
    & \multicolumn{3}{c}{Full \Macro{}}\\
    \cmidrule(lr){2-4}
    \cmidrule(l){5-7}
    Market
    & TR, \% & ASR & MDD, \%
    & TR, \% & ASR & MDD, \%\\
    \midrule
    BTC &  2.03 &  0.49 &  6.49
        &  3.03 &  0.61 &  5.41\\
    ETH & 12.73 &  1.26 & 20.87
        & 39.28 &  3.89 &  9.67\\
    DOT &  2.41 &  0.34 & 27.23
        & 13.79 &  0.97 & 15.89\\
    LTC & -8.66 & -0.58 & 22.03
        & 18.16 &  1.50 & 14.24\\
    \bottomrule
  \end{tabular}
\end{table*}

With memory, return and Sharpe ratio rise and MDD falls on all four assets. The return difference
is 26.55 percentage points on ETH and 26.82 on LTC. A general objection to architectural
complexity does not negate this positive result.

The table does not, however, separate the possible causes of the improvement. It provides no
distribution of the effect across repeated training runs, and the four assets' test periods
overlap substantially in calendar time. It therefore does not reveal whether the gain is robust
to random initialisation or to a change in market period.

In discussing Figure~6, the authors attribute the difference between the models to the no-memory
variant's late response to a fall in ETH. This illustrates the behaviour of two trained policies.
It cannot replace a comparison on a systematically selected sample of stress events, nor does it
show how much of the total improvement specifically comes from avoiding rare losses.

\subsubsection{Which problem does the mechanism solve?}

The motivation for memory combines several possible shortcomings. Rare events may receive too
few training updates, while Q-values may be unstable. In addition, an objective that maximises
mean reward may permit unacceptable risk. These problems require different solutions. A
``with-memory versus without-memory'' ablation cannot distinguish between the corresponding
explanations of the effect.

\paragraph{Insufficient use of rare transitions.}

If the problem is the distribution of training examples, their sampling frequency is what needs
to change. Prioritized Experience Replay raises the probability of resampling transitions with a
large TD error and uses importance-sampling corrections to control bias
\cite{schaul2016prioritized}. In a financial setting, one can also test event-based sampling:
identify sharp moves or large losses in the train set in advance and ensure their representation
in training batches.

PER does not guarantee special attention to financial disasters in particular: an ordinary noisy
transition may also have a large TD error. It does, however, directly test the hypothesis that
standard training allocates updates poorly across observations.

In \Macro{}, neither storing an entry nor retrieving it later depends directly on the size of a
loss or the rarity of the market event. FIFO favours recent observations, while retrieval favours
nearby states. A rare episode may be evicted from the table; while it remains there, its
contribution depends on how often it qualifies as a relevant neighbour.

This does not mean that all information about the episode disappears: it may remain in the
network parameters and the ordinary replay buffer. But the memory design provides no dedicated
mechanism for retaining crises over the long term or prioritising them.

Moreover, increasing the share of crisis examples without correcting their weights changes the
distribution against which the model is trained. If the purpose is to induce more cautious
behaviour, this change should be stated as part of the training objective, not merely as a remedy
for too few observations.

\paragraph{Stabilising and smoothing Q-values.}

The actual loss in~\eqref{eq:memory_loss} admits another explanation: memory creates a second
target by averaging stored estimates. This may reduce the variance of training targets and alter
the optimisation path even in ordinary, non-extreme states.

Algebraically, consider one update with fixed TD target $y$ and memory value $z=Q_M(s,a)$. The
isolated quadratic part of the loss is
\begin{equation}
  \ell(q)
  =
  (q-y)^2+\lambda_M(q-z)^2,
  \qquad q=Q_\theta(s,a).
  \label{eq:two_targets_loss}
\end{equation}
It is minimised at
\begin{equation}
  q_{\mathrm{eff}}
  =
  \frac{y+\lambda_M z}{1+\lambda_M}.
  \label{eq:effective_memory_target}
\end{equation}
Thus, the extra term shifts the target estimate towards the stored values. This is a property of
the two quadratic terms, not a solution of the full network-training problem.

Such smoothing is useful if neighbouring entries provide additional information about the same
conditional value. If nearby embeddings instead join states in which actions have different
consequences, averaging may obscure precisely the distinction that matters. This concern is
especially acute at a regime boundary, where the policy is expected to change its decision
quickly.

The initial observation that ``similar states have very different rewards'' does not identify the
source of the dispersion. It may reflect random noise, estimation error, or missing information
in the input state. Averaging may help in the first case; in the last, it cannot recover the
missing feature needed to distinguish the situations.

The stored values $v_j$ also contain the hyper-agent's own predictions. Memory provides no
independent check on its long-horizon estimates. It may smooth errors, but it may also propagate
a systematic error across similar states. Bootstrapped estimates are not themselves a defect:
ordinary Q-learning uses them too. Their presence does, however, limit the interpretation of the
table as a store of verified correct decisions.

Natural competitors for this hypothesis include other auxiliary Q-targets, value averaging, and
regularisation. A comparison with them would reveal whether neighbour retrieval and a separate
table are necessary, or whether the main effect comes from an additional stabilising term.

\paragraph{The policy's attitude to risk.}

Even perfect recall of every crisis transition does not compel an agent to avoid a risky
strategy. If the strategy has positive expected reward, maximising the expectation may favour it
over a safer alternative.

The \Macro{} loss includes an error that aligns Q-values, but no explicit constraint on tail
losses. The lower drawdown in the table is therefore an empirical training result, not a direct
consequence of an optimised risk criterion.

If the goal is to limit rare, large losses, it can be stated explicitly. For example, for total
reward $G_\pi$, consider
\begin{equation}
  J_{\mathrm{risk}}(\pi)
  =
  \mathbb E[G_\pi]
  -
  \lambda_R
  \operatorname{CVaR}_{\alpha}(-G_\pi),
  \label{eq:memory_risk_alternative}
\end{equation}
where $\lambda_R\geq0$, and $\operatorname{CVaR}_{\alpha}$ measures the mean loss in the upper
tail of the loss distribution; at $\alpha=0.95$, this is the worst 5\% of outcomes. Risk-sensitive
formulations using CVaR predate \Macro{}~\cite{chow2015cvar}.

Equation~\eqref{eq:memory_risk_alternative} defines an alternative problem; it does not describe
a hidden effect of memory. It requires comparison along the return--risk trade-off. Memory may
improve reward estimation, while a risk-sensitive objective changes the policy's preferences;
the two mechanisms can complement one another.

\subsubsection{What should be tested next}

The first step is to decide which explanation is being tested. Resampling data, smoothing
Q-targets, and changing the attitude to risk are not interchangeable ablations.

To test the rare-event explanation, ordinary replay, PER, event-based sampling, and \Macro{}'s
memory should be compared with the same low-level experts and hyper-agent. In addition to final
PnL, the experiment should measure how often crisis transitions take part in training, which
entries memory retrieves, and whether losses fall in new stress episodes. Without this evidence,
the rare-event explanation remains a conjecture.

Testing the role of retrieval requires a control with the same auxiliary loss term but a different
construction of $Q_M$. Nearest states can be compared with random entries of the same action and
similar age, as well as with uniform averaging over the retrieved neighbours. This separates the
selection of relevant transitions from the mere presence of a second training target. Q-value
quality should also be evaluated on held-out data: reducing error against the model's own
bootstrapped target does not yet establish a better reward prediction.

If the effect is attributed to risk management, the model should be compared with an explicit
risk objective. Tail losses, drawdowns, and foregone returns should be evaluated instead of
choosing a winner by total profit alone. Differences in market exposure should also be reported:
reducing losses simply by being in the market less often is not equivalent to better recognition
of dangerous states.

All comparisons should use the same training data, the same hyper-agent architecture, and
comparable budgets for tuning and updates. Training time and the cost of memory retrieval should
be reported separately. Repeated initialisations and non-overlapping test periods would help
distinguish a robust effect from the particulars of the selected trained policies.

The published ablation supports the usefulness of the added module in this particular system,
but does not establish which training defect it remedies. Justifying memory as a mechanism for
handling extreme events requires precisely this separation of mechanisms. Otherwise, the high
return confirms a successful combination of components, but does not explain why episodic memory
should be preferred to existing ways of using experience and managing risk in the same task.

% ------------------------------------------------------------
% References
% ------------------------------------------------------------

\begin{thebibliography}{99}
  \setlength{\itemsep}{0.55ex}

  \bibitem{zong2024macrohft}
  C.~Zong, C.~Wang, M.~Qin, L.~Feng,
  X.~Wang, and B.~An.
  \newblock
  MacroHFT: Memory Augmented Context-aware
  Reinforcement Learning On High Frequency Trading.
  \newblock
  \emph{Proceedings of the 30th ACM SIGKDD
  Conference on Knowledge Discovery and Data Mining},
  2024, pp.~4712--4721.
  \newblock
  DOI:
  \href{https://doi.org/10.1145/3637528.3672064}
  {10.1145/3637528.3672064}.

  \bibitem{qin2024earnhft}
  M.~Qin, S.~Sun, W.~Zhang, H.~Xia,
  X.~Wang, and B.~An.
  \newblock
  EarnHFT: Efficient Hierarchical Reinforcement
  Learning for High Frequency Trading.
  \newblock
  \emph{Proceedings of the AAAI Conference
  on Artificial Intelligence},
  38(13), 2024, pp.~14669--14676.
  \newblock
  DOI:
  \href{https://doi.org/10.1609/aaai.v38i13.29384}
  {10.1609/aaai.v38i13.29384}.

  \bibitem{deng2016deep}
  Y.~Deng, F.~Bao, Y.~Kong, Z.~Ren,
  and Q.~Dai.
  \newblock
  Deep Direct Reinforcement Learning for
  Financial Signal Representation and Trading.
  \newblock
  \emph{IEEE Transactions on Neural Networks
  and Learning Systems},
  28(3), 2017, pp.~653--664
  (published online in 2016).
  \newblock
  DOI:
  \href{https://doi.org/10.1109/TNNLS.2016.2522401}
  {10.1109/TNNLS.2016.2522401}.

  \bibitem{theate2021trading}
  T.~Th\'eate and D.~Ernst.
  \newblock
  An Application of Deep Reinforcement
  Learning to Algorithmic Trading.
  \newblock
  \emph{Expert Systems with Applications},
  173, 2021, article 114632.
  \newblock
  DOI:
  \href{https://doi.org/10.1016/j.eswa.2021.114632}
  {10.1016/j.eswa.2021.114632}.

  \bibitem{zhu2022perturbation}
  T.~Zhu and W.~Zhu.
  \newblock
  Quantitative Trading through Random
  Perturbation Q-Network with Nonlinear
  Transaction Costs.
  \newblock
  \emph{Stats},
  5(2), 2022, pp.~546--560.
  \newblock
  DOI:
  \href{https://doi.org/10.3390/stats5020033}
  {10.3390/stats5020033}.

  \bibitem{almgren2001optimal}
  R.~Almgren and N.~Chriss.
  \newblock
  Optimal Execution of Portfolio Transactions.
  \newblock
  \emph{The Journal of Risk}, 3(2), 2001,
  pp.~5--39.
  \newblock
  DOI:
  \href{https://doi.org/10.21314/JOR.2001.041}
  {10.21314/JOR.2001.041}.

  \bibitem{byrd2020abides}
  D.~Byrd, M.~Hybinette, and T.~H.~Balch.
  \newblock
  ABIDES: Towards High-Fidelity Multi-Agent
  Market Simulation.
  \newblock
  \emph{Proceedings of the ACM SIGSIM Conference
  on Principles of Advanced Discrete Simulation},
  2020, pp.~11--22.
  \newblock
  DOI:
  \href{https://doi.org/10.1145/3384441.3395986}
  {10.1145/3384441.3395986}.

  \bibitem{lin2018episodic}
  Z.~Lin, T.~Zhao, G.~Yang, and L.~Zhang.
  \newblock
  Episodic Memory Deep Q-Networks.
  \newblock
  arXiv:1805.07603, 2018.
  \newblock
  \url{https://arxiv.org/abs/1805.07603}.

  \bibitem{pritzel2017neural}
  A.~Pritzel, B.~Uria, S.~Srinivasan,
  A.~Puigdom\`enech Badia, O.~Vinyals,
  D.~Hassabis, D.~Wierstra, and C.~Blundell.
  \newblock
  Neural Episodic Control.
  \newblock
  \emph{Proceedings of the 34th International
  Conference on Machine Learning},
  PMLR 70, 2017, pp.~2827--2836.
  \newblock
  \url{https://proceedings.mlr.press/v70/pritzel17a.html}.

  \bibitem{schaul2016prioritized}
  T.~Schaul, J.~Quan, I.~Antonoglou,
  and D.~Silver.
  \newblock
  Prioritized Experience Replay.
  \newblock
  \emph{International Conference
  on Learning Representations}, 2016.
  \newblock
  \url{https://arxiv.org/abs/1511.05952}.

  \bibitem{chow2015cvar}
  Y.~Chow, A.~Tamar, S.~Mannor,
  and M.~Pavone.
  \newblock
  Risk-Sensitive and Robust Decision-Making:
  a CVaR Optimization Approach.
  \newblock
  \emph{Advances in Neural Information
  Processing Systems}, 28, 2015,
  pp.~1522--1530.
  \newblock
  \url{https://arxiv.org/abs/1506.02188}.

  \bibitem{macrohftcode}
  ZONG0004.
  \newblock
  MacroHFT: official implementation.
  \newblock
  Version reviewed:
  commit
  \texttt{31e5ef41f93b2aea6e63e6ae2267675c997a8e54}.
  Files
  \path{RL/util/memory.py}
  and \path{RL/agent/high_level.py}.
  \newblock
  \url{https://github.com/ZONG0004/MacroHFT/tree/31e5ef41f93b2aea6e63e6ae2267675c997a8e54}.


\end{thebibliography}

\end{document}
